% Generated by Listas/make_statements.py from pratica-lecture-02.md, item 2. Do not edit.
\textbf{Source.} Tadelis (2013), \emph{Game Theory: An Introduction}, Exercise 5.6 ``Hawk-Dove'', p.~96 {[}PDF p.~115{]}. Answer in the source: none (the preface, p.~xiii, says about half of the exercises are solved on Princeton UP's website). Reproduced verbatim (with the book's footnote 12).



\begin{quote}

\textbf{Hawk-Dove:} The following game has been widely used in evolutionary biology to understand how fighting and display strategies by animals could coexist in a population. For a typical Hawk-Dove game there are resources to be gained (e.g., food, mates, territories), denoted as \(v\). Each of two players can choose to be aggressive, as Hawk (\(H\)), or compromising, as Dove (\(D\)). If both players choose \(H\) then they split the resources but lose some payoff from injuries, denoted as \(k\). Assume that \(k > \frac{v}{2}\). If both choose \(D\) then they split the resources but engage in some display of power that carries a display cost \(d\), with \(d < \frac{v}{2}\). Finally, if player \(i\) chooses \(H\) while \(j\) chooses \(D\) then \(i\) gets all the resources while \(j\) leaves with no benefits and no costs.



\begin{enumerate}

\def\labelenumi{\alph{enumi}.}

\item

  Describe this game in a matrix.

\item

  Assume that \(v = 10\), \(k = 6\), and \(d = 4\). What outcomes can be supported as pure-strategy Nash equilibria?\(^{12}\)

\end{enumerate}



\emph{Footnote 12.} In the evolutionary biology literature, the analysis performed is of a very different nature. Instead of considering the Nash equilibrium analysis of a static game, the analysis is a dynamic one in which successful strategies ``replicate'' in a large population. This analysis is part of a methodology called evolutionary game theory. For more on the subject see Gintis (2000).

\end{quote}

